% Pursuing Stacks (cotes 134-2 to 134-8): archive cote 134-2, site page 50 (PDF page 51).
% appendix (letters to Breen). Draft compared: ps1.tex:1250-1261.
% AI-transcribed by Claude Opus 5.5 (claude-opus-5-5, adaptive thinking, effort high), Message Batches API, batch ps1 (msgbatch_01Rv5hCaNTtUt4JHx1yqoHPL), 2026-11-12; not yet checked by eye.
\scanpage{134-2/50}
\mspage{\hand{35}} \mspage{3}

``(n-1)-systèmes locaux sur X'' i.e. des (n-1)- champs localement constants
sur X, et par là élucider complètement l'énonc\over{e}{é}\deleted{\unsure{s}} hypothétique (B) ci-dessus,
il faudrait donc expliciter comment un n-groupoide \hand{$\mathcal{C}$}$_n$ se récupère, à
n-équivalence près, par la connaissance de la n-catégorie
\hand{$\mathcal{T}$}$_n$ = (n-\underline{Hom}) (\hand{$\mathcal{C}$}$_n$, ( (n-1)-Cat)) des (n-1)- systèmes locaux sur \hand{$\mathcal{C}$}$_n$. On
aurait envie de dire que \hand{$\mathcal{C}$}$_n$ est la catégorie des ``n-foncteurs fibres'' sur
\hand{$\mathcal{T}$}$_n$, i.e. des n- foncteurs \hand{$\mathcal{T}$}$_n$ \hand{$\longrightarrow$}((n-1)-Cat) ayant certaines proprié-%
tés d'exactitude (pour n = 1, c'était la condition d'être les foncteurs
image inverse pour un morphisme de topos, i.e. de commuter aux $\underset{\hand{$\longrightarrow$}}{\mathrm{lim}}$ quelcon-%
ques et aux $\underset{\hand{$\longleftarrow$}}{\mathrm{lim}}$ finies ...) C'est ici que se matérialise la peur, ex-%
primée dans ma\deleted{\unsure{s}} précédente lettre, qu' on finisse par tomber sur la notion
de n-topos et morphismes de tels ! \hand{$\mathcal{T}$}$_n$ serait un n-topos\hand{,} ( appelé
le ``n-topos classifiant du n-groupoide \hand{$\mathcal{C}$}$_n$), ((n-1)-Cat serait le
n-topos ``ponctuel'' type , et \hand{$\mathcal{C}$}$_n$ s'interprète mod n-équivalence
comme la n-catégorie \inserted{des} ``n-points'' du n-topos classifiant \hand{$\mathcal{T}$}$_n$. Brr !

Si on espère encore pouvoir définir un bon vieux (1)-topos classifiant
pour un n-groupoide \hand{$\mathcal{C}$}$_n$, comme solution d'un Pb universel, je ne vois
guère que le Pb universel suivant: pour tout topos T, considérons
\margin{\hand{cela} sem- \hand{ble} sus- \hand{pect!}}\underline{Hom} (\hand{$\Pi$}$_n$(T), \hand{$\mathcal{C}$}$_n$). C'est une n-catégorie, mais prenons-en la 1-catégorie
tronquée \hand{$\tau$}$_1$\underline{Hom} (\hand{$\Pi$}$_n$(T), \hand{$\mathcal{C}$}$_n$). Pour T variable, on\hand{/}voudrait 2-représenter
le 2-foncteur contravariant (Top)$^{o}$ \hand{$\longrightarrow$} (1-Cat) par un topos classi-%
fiant \hand{$\mathcal{B}$} = \hand{$\mathcal{B}_{\mathcal{C}_n}$} , donc trouver un \hand{$\Pi$}$_n$(\hand{$\mathcal{B}$}) \hand{$\xrightarrow{\varphi}$} \hand{$\mathcal{C}$}$_n$ 2-univer-%
sel en le sens que pour tout T, le foncteur
\[ \underline{\mathrm{Hom}}_{(\mathrm{Top})}(T , \hand{$\mathcal{B}$} ) \hand{$\xrightarrow{u \mapsto \varphi \circ \Pi_n(u)}$} \hand{$\tau$}_1\underline{\mathrm{Hom}} (\hand{$\Pi$}_n(T), \hand{$\mathcal{C}$}_n) \]
soit une équivalence. Pour n=1 on sait que le topos classifiant de \hand{$\mathcal{C}$}$_1$
au sens usuel fait l'affaire, mais pour n=2 déja, je \over{\unsure{d}}{d}oute que ce Pb
universel ait une solution. C'est peut-être lié au fait que le
``théorème de van \over{\unsure{K}}{K}ampen'', qu'on peut exprimer en disant que le 2-foncteur
T \hand{$\longrightarrow\Pi$}$_1$(T) des topos localement 1-connexes vers les groupoides trans-%
forme (à 1-équivalence près) sommes amalgamées en sommes amalgamées ( et
plus \unsure{générahement|généralement} commute aux 2-limites inductives), n'est sans doute plus

\margin{50}
